\chapter{Architektura systemu UniVerify}
\label{ch:architektura}

Rozdział przedstawia budowę systemu UniVerify w~postaci, na której opierają się
kolejne rozdziały pomiarowe i~projektowe. Opisuje model danych, kontrakt rejestru,
pakiety współdzielone oraz aplikację webową, a~na końcu zbiera je w~trzech
przepływach end-to-end: wydania, weryfikacji i~rewokacji. Pomiary kosztów pozostają
poza zakresem tego rozdziału --- omawiają je rozdziały~\ref{ch:gaz} i~\ref{ch:l2}.

\section{Model danych}
\label{sec:model-danych}

Podstawową jednostką zapisu on-chain nie jest pojedynczy dyplom, lecz \emph{partia}
(\emph{batch}). Wystawca grupuje dyplomy w~drzewo Merkle i~publikuje w~kontrakcie
jedynie jego korzeń. Dzięki temu koszt zapisu nie zależy od liczby dokumentów
w~partii --- jego uzasadnienie ilościowe podaje rozdział~\ref{ch:gaz}. Trójka
identyfikująca partię to:

\begin{itemize}
  \item \code{issuer} --- adres Ethereum portfela autoryzowanego przez uczelnię;
  \item \code{batchId} --- 64-bitowy numer kolejny partii, unikalny w~obrębie
        wystawcy;
  \item \code{merkleRoot} --- 32-bajtowy korzeń drzewa Merkle danej partii.
\end{itemize}

Pojedynczy dyplom reprezentowany jest przez \emph{kopertę}
(\emph{envelope}) złożoną z~dwóch części: \code{payload} --- danych dyplomu
(identyfikator uczelni, dane studenta, stopień, data wydania, numer dyplomu) ---
oraz \code{proof} typu \code{MERKLE\_BATCH}, zawierającego ścieżkę Merkle wiążącą
liść dokumentu z~korzeniem partii. Weryfikacja polega na odtworzeniu korzenia
z~liścia i~ścieżki oraz porównaniu go z~korzeniem odczytanym z~kontraktu.

Status dyplomu opisuje wartość \code{StatusCode}: $0$~---~nieznany
(\emph{Unknown}), $1$~---~ważny (\emph{Valid}), $2$~---~unieważniony
(\emph{Revoked}). Status nieznany oznacza, że liść nie należy do żadnej znanej
partii; unieważnienie omówiono w~sekcji~\ref{sec:przeplywy}.

\section{Kontrakt \code{DiplomaRegistry}}
\label{sec:kontrakt}

Kontrakt \code{DiplomaRegistry} udostępnia funkcje na trzech poziomach uprawnień.

\begin{description}
  \item[Właściciel (\emph{owner}).] Zarządza uczelniami i~wystawcami: rejestruje
        uczelnię wraz z~metadanymi (nazwa, kraj, strona, identyfikator
        akredytacji), przypisuje wystawcom uczelnie oraz odbiera im uprawnienia.
        Właściciel korzysta z~dwustopniowego przekazania własności
        (\code{transferOwnership} $\rightarrow$ \code{acceptOwnership}), co
        rozdział~\ref{ch:governance} wykorzystuje do zainstalowania modelu
        zarządzania bez zmiany kontraktu.
  \item[Wystawca (\emph{issuer}).] Publikuje korzenie partii
        (\code{issueBatchRoot}) oraz unieważnia pojedyncze dyplomy
        (\code{revokeFromBatch}), przy czym unieważnienie wymaga podania dowodu
        Merkle potwierdzającego przynależność dokumentu do partii.
  \item[Publiczny odczyt (\emph{view}).] Funkcje bezstanowe dostępne dla każdego:
        \code{statusWithProof} (status wraz z~weryfikacją dowodu), \code{getBatch}
        (korzeń partii), \code{isIssuer} oraz \code{issuerUniversityId}.
\end{description}

Kontrakt utrzymuje metadane uczelni, mapowanie wystawca~$\rightarrow$~uczelnia,
korzenie Merkle indeksowane parą \code{(issuer, batchId)} oraz zbiór unieważnień.
Status uczelni opisuje osobne wyliczenie: nieznana~(0), aktywna~(1),
zawieszona~(2), unieważniona~(3).

\section{Pakiety współdzielone}
\label{sec:pakiety}

System zorganizowano jako monorepozytorium, w~którym logika kryptograficzna żyje
w~jednym miejscu i~jest współdzielona przez aplikację webową, CLI oraz SDK
(rys.~\ref{fig:architektura}).

\begin{figure}[H]
\centering
\begin{tikzpicture}[node distance=6mm and 8mm]
  % Warstwa klientów
  \node[blok] (web) {aplikacja webowa\\\code{Next.js}};
  \node[blok, right=of web] (cli) {CLI\\\code{issue / verify / gov}};
  \node[blok, right=of cli] (sdk) {SDK\\\code{UniverifySdk}};
  % Rdzeń
  \node[blokakcent, below=9mm of cli, minimum width=68mm]
    (core) {\code{verifier-core}\\haszowanie, drzewa Merkle, schemat ładunku, ABI};
  % On-chain
  \node[blokszary, below=9mm of core, minimum width=68mm]
    (chain) {kontrakt \code{DiplomaRegistry} (L1 / L2)\\korzenie partii, uprawnienia wystawców, unieważnienia};
  \draw[strz] (web) -- (core);
  \draw[strz] (cli) -- (core);
  \draw[strz] (sdk) -- (core);
  \draw[strz] (core) -- node[etyk, right=1mm] {odczyt / transakcje RPC} (chain);
\end{tikzpicture}
\caption[Architektura systemu UniVerify]{Architektura systemu UniVerify: trzy interfejsy klienckie współdzielą
jedną implementację kryptograficzną (\code{verifier-core}), a~na łańcuchu
utrzymywany jest wyłącznie kontrakt rejestru.}
\label{fig:architektura}
\end{figure}

\begin{description}
  \item[\code{verifier-core}.] Warstwa fundamentalna. Eksportuje
        \code{hashPayload} (kanonizacja + Keccak256), \code{computeMerkleLeaf}
        (haszowanie liścia z~separacją domen), \code{buildMerkleFromLeaves}
        (budowa drzewa i~zwrot korzenia oraz dowodów), \code{computeRootFromProof}
        (odtworzenie korzenia z~liścia i~dowodu), a~także ABI kontraktu i~schemat
        walidacji ładunku (Zod). Musi być zbudowana przed pozostałymi pakietami.
  \item[\code{sdk}.] Klasa \code{UniverifySdk} opakowuje odczyty z~kontraktu
        i~pełną weryfikację: \code{verify} (weryfikacja koperty wraz z~metadanymi
        uczelni), \code{status}, \code{getBatch}, \code{getIssuer},
        \code{getUniversity} oraz odczyt informacji o~własności.
  \item[\code{verifier-cli}.] Skrypty wiersza poleceń: \code{issue.ts}
        i~\code{verify.ts}, a~także grupa poleceń \code{gov} obsługująca model
        zarządzania z~rozdziału~\ref{ch:governance}.
\end{description}

Rozdzielenie logiki do \code{verifier-core} sprawia, że ta sama, przetestowana
implementacja haszowania i~drzew Merkle jest używana zarówno przy wystawianiu, jak
i~przy weryfikacji --- eliminuje to ryzyko rozjazdu między stroną wydającą
a~sprawdzającą.

\section{Aplikacja webowa}
\label{sec:web}

Interfejs użytkownika to aplikacja Next.js udostępniająca trasy dopasowane do ról
w~systemie (tab.~\ref{tab:web-routes}) oraz dwa punkty API.

\begin{table}[H]
\centering
\caption{Trasy aplikacji webowej i~ich role.}
\label{tab:web-routes}
\begin{tabular}{@{}lll@{}}
\toprule
Trasa & Rola & Przeznaczenie \\
\midrule
\code{/issuer}   & wystawca  & budowa partii Merkle, publikacja korzeni, eksport dowodów \\
\code{/verifier} & publiczna & weryfikacja koperty i~sprawdzenie statusu on-chain \\
\code{/revoke}   & wystawca  & unieważnianie dyplomów przez dowód Merkle \\
\code{/present}  & posiadacz & selektywne ujawnianie pól (rozdz.~\ref{ch:prywatnosc}) \\
\code{/admin}    & właściciel & zarządzanie wystawcami i~uczelniami (rozdz.~\ref{ch:governance}) \\
\bottomrule
\end{tabular}
\end{table}

Punkty \code{/api/verify} (pełna weryfikacja z~metadanymi uczelni) oraz
\code{/api/status} (lekki odczyt statusu) chroni wspólny \emph{guard}: opcjonalne
uwierzytelnianie kluczem API (wymuszane tylko, gdy ustawiono zmienną środowiskową
z~kluczami; akceptuje nagłówek \code{Bearer} lub \code{X-API-Key}) oraz ograniczanie
liczby żądań (30~żądań na~60~s na~adres IP, w~przesuwnym oknie w~pamięci).

\section{Przepływy: wydanie, weryfikacja, rewokacja}
\label{sec:przeplywy}

Trzy główne przepływy spinają opisane elementy w~całość
(rys.~\ref{fig:przeplywy}).

\begin{figure}[H]
\centering
\begin{tikzpicture}[node distance=14mm and 24mm]
  % Wydanie
  \node[blok] (uczelnia) {uczelnia\\(wystawca)};
  \node[blokakcent, right=of uczelnia] (drzewo) {drzewo Merkle\\(off-chain, bez gazu)};
  \node[blokszary, right=of drzewo] (kontrakt) {kontrakt\\\code{DiplomaRegistry}};
  \node[blok, below=of drzewo] (absolwent) {absolwent\\(koperta: \code{payload} + dowód)};
  \node[blok, below=of kontrakt] (weryfikator) {weryfikator};
  \draw[strz] (uczelnia) -- node[etyk, above=1mm] {liście\\dyplomów} (drzewo);
  \draw[strz] (drzewo) -- node[etyk, above=1mm] {\code{issueBatchRoot}\\(korzeń)} (kontrakt);
  \draw[strz] (drzewo) -- node[etyk, left=1mm] {eksport kopert} (absolwent);
  \draw[strz] (absolwent) -- node[etyk, above=1mm] {koperta} (weryfikator);
  \draw[strz] (weryfikator) -- node[etyk, right=1mm, align=left]
    {\code{statusWithProof}:\\porównanie korzeni} (kontrakt);
  \draw[strz, dashed] (uczelnia.south) -- ++(0,-32mm) -|
    node[etyk, below, pos=0.25] {\code{revokeFromBatch} (dowód Merkle)}
    (kontrakt.south east);
\end{tikzpicture}
\caption[Przepływy: wydanie, weryfikacja, rewokacja]{Przepływy w~systemie: wydanie (uczelnia publikuje wyłącznie korzeń
partii), weryfikacja (odtworzenie korzenia z~koperty i~porównanie ze stanem
kontraktu) oraz rewokacja (linia przerywana).}
\label{fig:przeplywy}
\end{figure}

\paragraph{Wydanie.} Wystawca buduje off-chain drzewo Merkle z~liści dyplomów
(\code{computeMerkleLeaf} + \code{buildMerkleFromLeaves}), publikuje w~kontrakcie
sam korzeń (\code{issueBatchRoot}) i~eksportuje każdemu absolwentowi kopertę zawierającą
jego \code{payload} oraz dowód Merkle. Cała praca kryptograficzna poza jedną
transakcją odbywa się bez gazu.

\paragraph{Weryfikacja.} Weryfikator (przez \code{/verifier}, API lub SDK)
odtwarza z~koperty liść, z~liścia i~dowodu --- korzeń, po czym porównuje go
z~korzeniem odczytanym z~kontraktu dla pary \code{(issuer, batchId)}. Pozytywny
wynik oznacza, że dokument należy do partii opublikowanej przez autoryzowanego
wystawcę i~nie został unieważniony.

\paragraph{Rewokacja.} Aby unieważnić pojedynczy dyplom, wystawca wywołuje
\code{revokeFromBatch}, podając dowód Merkle dokumentu. Kontrakt dodaje liść do
zbioru unieważnień, przez co kolejne odczyty statusu zwracają \code{Revoked}, mimo
że korzeń partii pozostaje niezmieniony. Pozwala to cofnąć pojedynczy dokument bez
naruszania pozostałych dyplomów tej samej partii.
